\chapter{Transformers and Attention Mechanisms From Scratch}
\label{chap:ch10}

\section{The Fundamental Limitation of Convolutional Networks}
In a standard convolutional layer, each output neuron only sees a small local patch (e.g., $3 \times 3$ pixels). To capture relationships between two pixels separated by 80 pixels, a CNN must stack dozens of layers so that their receptive fields gradually overlap.

In satellite cloud removal, this is a severe weakness:
\begin{itemize}
    \item A single cloud can cover a $60 \times 60$ pixel area.
    \item The terrain on the southern edge of the cloud might be an olive orchard.
    \item An identical olive orchard may be clearly visible 70 pixels away on the northern edge of the cloud!
    \item A convolutional network struggles to connect these distant parcels because the information must diffuse slowly across dozens of intermediate layers.
\end{itemize}

\section{The Attention Mechanism: Queries, Keys, and Values}
Vaswani et al.~\cite{vaswani2017attention} introduced the \textbf{Attention Mechanism} to allow any element in a sequence or image to interact directly with any other element in a single computational step.

\begin{conceptbox}{Everyday Analogy: A Filing Cabinet Search}
Imagine searching for a recipe in a kitchen library:
\begin{itemize}
    \item \textbf{Query ($Q$):} Your search request (e.g., "warm vegetable stew").
    \item \textbf{Key ($K$):} The label or title written on each folder tab in the cabinet (e.g., "Folder 1: Chocolate Cakes", "Folder 2: Autumn Soups and Stews", "Folder 3: Grilled Steaks").
    \item \textbf{Attention Score ($Q \cdot K^T$):} You compare your search query against every folder label. Folder 2 receives a high similarity score ($0.92$), while Folders 1 and 3 receive near zero.
    \item \textbf{Value ($V$):} The actual contents (recipes) inside the folders.
    \item \textbf{Output:} You take a weighted sum of the contents, extracting 92\% of your information from Folder 2 and almost nothing from the others!
\end{itemize}
\end{conceptbox}

\section{The Scaled Dot-Product Attention Equation}
Formally, given Query matrix $\mathbf{Q}$, Key matrix $\mathbf{K}$, and Value matrix $\mathbf{V}$:
\begin{equation}
\text{Attention}(\mathbf{Q}, \mathbf{K}, \mathbf{V}) = \text{Softmax}\left(\frac{\mathbf{Q} \mathbf{K}^T}{\sqrt{d_k}}\right) \mathbf{V}
\end{equation}
\begin{itemize}
    \item $\mathbf{Q} \mathbf{K}^T$: Computes pairwise dot-product similarities between every query and key.
    \item $\frac{1}{\sqrt{d_k}}$: Scaling factor based on feature dimension $d_k$ to prevent the dot products from growing excessively large for large dimensions (which would push softmax into near-zero gradient regions).
    \item $\text{Softmax}(\cdot)$: Normalizes similarity scores along each row so that they sum to $1.0$, turning them into valid attention probabilities.
    \item Multiply by $\mathbf{V}$: Assembles the final output as a convex combination of value representations.
\end{itemize}

\begin{mathbox}{A Complete Numerical Attention Calculation}
Let us trace an exact numerical calculation for 2 tokens (dimension $d_k = 2$):
\[
\mathbf{Q} = \begin{bmatrix} 1.0 & 0.0 \\ 0.0 & 2.0 \end{bmatrix}, \quad
\mathbf{K} = \begin{bmatrix} 1.0 & 1.0 \\ 0.0 & 1.0 \end{bmatrix}, \quad
\mathbf{V} = \begin{bmatrix} 10.0 & 20.0 \\ 30.0 & 40.0 \end{bmatrix}
\]
Here $d_k = 2$, so $\sqrt{d_k} = \sqrt{2} \approx 1.414$.

\textbf{Step 1: Compute Raw Attention Scores ($\mathbf{Q}\mathbf{K}^T$)}
\[
\mathbf{Q}\mathbf{K}^T = \begin{bmatrix} 1.0 & 0.0 \\ 0.0 & 2.0 \end{bmatrix} \begin{bmatrix} 1.0 & 0.0 \\ 1.0 & 1.0 \end{bmatrix} = \begin{bmatrix} 1.0 & 0.0 \\ 2.0 & 2.0 \end{bmatrix}
\]

\textbf{Step 2: Scale by $\frac{1}{\sqrt{d_k}}$}
\[
\mathbf{S} = \frac{\mathbf{Q}\mathbf{K}^T}{1.414} = \begin{bmatrix} 0.707 & 0.000 \\ 1.414 & 1.414 \end{bmatrix}
\]

\textbf{Step 3: Apply Row-Wise Softmax ($\mathbf{A} = \text{Softmax}(\mathbf{S})$)}
For Row 0:
\[
e^{0.707} \approx 2.028, \quad e^{0.000} = 1.000 \implies \text{Sum} = 3.028
\]
\[
\mathbf{A}_{0,0} = \frac{2.028}{3.028} \approx \mathbf{0.670}, \quad \mathbf{A}_{0,1} = \frac{1.000}{3.028} \approx \mathbf{0.330}
\]
For Row 1 (both values equal $1.414$, so equal probabilities):
\[
\mathbf{A}_{1,0} = \mathbf{0.500}, \quad \mathbf{A}_{1,1} = \mathbf{0.500}
\]
Attention weight matrix:
\[
\mathbf{A} = \begin{bmatrix} 0.670 & 0.330 \\ 0.500 & 0.500 \end{bmatrix}
\]

\textbf{Step 4: Multiply by Values ($\mathbf{Y} = \mathbf{A}\mathbf{V}$)}
\begin{align*}
\mathbf{Y}_{0} &= 0.670 \begin{bmatrix} 10.0 & 20.0 \end{bmatrix} + 0.330 \begin{bmatrix} 30.0 & 40.0 \end{bmatrix} = \begin{bmatrix} 6.70 + 9.90 & 13.40 + 13.20 \end{bmatrix} = \begin{bmatrix} \mathbf{16.60} & \mathbf{26.60} \end{bmatrix} \\
\mathbf{Y}_{1} &= 0.500 \begin{bmatrix} 10.0 & 20.0 \end{bmatrix} + 0.500 \begin{bmatrix} 30.0 & 40.0 \end{bmatrix} = \begin{bmatrix} 5.0 + 15.0 & 10.0 + 20.0 \end{bmatrix} = \begin{bmatrix} \mathbf{20.00} & \mathbf{30.00} \end{bmatrix}
\end{align*}
The output vectors synthesize information from both inputs weighted by their contextual relevance!
\end{mathbox}

\section{The Complexity Problem: Spatial Attention vs. Channel Attention}
For an image patch of size $128 \times 128$, the total number of spatial pixels is:
\[
N = H \times W = 128 \times 128 = 16{,}384\text{ pixels}
\]
Standard spatial self-attention computes an affinity matrix between every pixel and every other pixel:
\[
\text{Affinity Matrix Size} = N \times N = 16{,}384 \times 16{,}384 = \mathbf{268{,}435{,}456\text{ elements!}}
\]
At 4 bytes per float, storing this single attention matrix requires \textbf{1.07 GB of VRAM per attention head per batch item!} A 4-stage network would crash from out-of-memory errors instantly!

\subsection{The Restormer Breakthrough: Transposed Attention (MDTA)}
In Chapter 11, we will see how ECRformer solves this: instead of computing attention across the spatial pixels ($N \times N$), it computes \textbf{Transposed Attention across the feature channels} ($C \times C$, Zamir et al.~\cite{zamir2022restormer}, El-Nouby et al.~\cite{elnouby2021xcit}). Because $C \in \{32, 64, 128, 256\}$, a $256 \times 256$ matrix has only $65{,}536$ elements ($<1\text{ MB}$), reducing computational complexity from $\mathcal{O}((HW)^2)$ to linear in image area: $\mathcal{O}(C^2 HW)$!

\section{Summary of Chapter 10}
\begin{enumerate}
    \item Convolutional receptive fields grow slowly, making them inefficient at modeling distant contextual dependencies across large clouds.
    \item Attention computes dynamic similarity scores between queries ($Q$) and keys ($K$) to aggregate relevant values ($V$).
    \item Scaled dot-product attention divides by $\sqrt{d_k}$ to prevent gradient saturation.
    \item Conventional spatial self-attention scales quadratically ($\mathcal{O}((HW)^2)$), making it computationally infeasible for high-resolution satellite imagery.
    \item Transposed cross-covariance attention computes attention across channels ($C \times C$), achieving linear spatial complexity $\mathcal{O}(C^2 HW)$.
\end{enumerate}

\section{Practice Exercises}
\begin{enumerate}
    \item Why is the scaling factor $\frac{1}{\sqrt{d_k}}$ necessary before applying the Softmax operation?
    \item Calculate the memory required to store an attention matrix for a $256 \times 256$ satellite image under standard spatial self-attention.
    \item If two key vectors $\mathbf{K}_1$ and $\mathbf{K}_2$ have identical dot products with query $\mathbf{Q}$, what will be their relative attention weights after Softmax?
\end{enumerate}

\subsection*{Solutions}
\begin{enumerate}
    \item For large feature dimensions $d_k$, the dot product $\sum_{i=1}^{d_k} q_i k_i$ sums many terms, growing in variance proportional to $d_k$. Without division by $\sqrt{d_k}$, extreme values push the Softmax function into saturation regions where its derivative is near zero, causing vanishing gradients.
    \item Spatial pixels: $N = 256 \times 256 = 65{,}536$. Attention matrix: $65{,}536 \times 65{,}536 = 4{,}294{,}967{,}296$ floats. At 4 bytes/float, this requires $17.18\text{ GB}$ of GPU memory for a single matrix!
    \item Because the Softmax exponentiates equal numerical scores ($e^s = e^s$), both keys will receive exactly identical attention weights ($50\%$ each if there are only two keys).
\end{enumerate}
